% s-152-p152-m002; original PDF 152, 153
\Aphorism{s-152-p152-m002}{240}
Es giebt Minima im „Pathos der Distanz“, Stunden, wo die beständige Spannung nach aussen und unten nachlässt und die Seele sich nach den überwundenen Entwickelungsstufen zurücksehnt, wo man Kind, Volk, Thier, Element zu werden verlangt. Zu diesen Anwandlungen von Seelenmüdigkeit zählt die Liebe, als eine sublime Begierde nach Heimath, Mutterschooss, Ruhe in der Gattung; in ihr ist, oft als Wille zum Ende verkleidet, ein unstillbarer Wille zum Anfang, zu den Ursprüngen unseres Daseins, der Wunsch unterzutauchen und zurückzusinken dorthin, woher wir gekommen. Dem Manne, der es eines Tages schrecklich empfindet, sich von Natur und Instincten gelöst zu haben, geht im Weibe eine Möglichkeit auf, zur Natur zurückzuflüchten und seinen Frieden mit dem Leben zu machen. Nun, dies hat seine Zeit, und an einem neuen Morgen wird der neue Krieg erklärt: schlimm, wenn jener Friede mehr sein wollte als ein schöner Zwischenfall.
\par

